\clearpage\begingroup\SourceLeaf{40}\HistoricalHead{35}{R}{GAUSS AN GERLING.}
\begin{center}
[2.]\\[.7em]
\textsc{Gauss} an \textsc{Gerling}. Göttingen, 17. Januar 1852.
\end{center}
\noindent .\ .\ .\ .\ .\ .\ .\ .\ .\ .\ .\ .\ .\ .\ .\ .\ .\ .\ .\ .\ .
\par\noindent\hspace*{\parindent}Es würde mich sehr freuen, wenn Sie Ihre Pendelschwingungsbeobb.
irgendwo ausführlich bekannt machen und dann zugleich eine historisch literarische
Zusammenstellung von allem \textls[120]{erheblichen}, was bisher darüber gearbeitet
ist, geben wollten. Ich habe nur hie und da einige dürftige Notizen
gesehen, bin also in der Sache nur wenig bekannt, und scheue die Mühe des
Zusammensuchens, oder vielmehr, es ist mir unmöglich, auf solches Zusammensuchen
Zeit zu verwenden. Es wollte mir scheinen, dass jede theoretische
Behandlung durchaus unbefriedigend sein müsse, wo nicht \textls[120]{der Hergang an
der Aufhängungsstelle} gründlich mit berücksichtigt wird, was nicht leicht
sein möchte, ausgenommen den einfachsten Fall, wo das Pendel oben mit
einer conischen Spitze auf einer harten Unterlage aufsitzt. Eine solche Spitze
ist aber im streng mathem[atischen] Sinn durch die Kunst nicht herzustellen.
Nehmen Sie zunächst anstatt der Spitze eine prismatische Schneide, so ist kein
Zweifel, dass wenn dieselbe eine gewisse Länge hat, die Schwingungen unverändert
in Einer unveränderl[ichen] Azimutalebene geschehen werden. Fragt
sich nur, wenn man die Länge der Schneide immer kürzer nimmt, bei welcher
Länge derselben die Rotation der Erde die Wirkung haben wird, dass diese
Schneide an einem Ende sich in die Höhe hebt, so dass das Pendel eine
Zeitlang bloss auf dem andern Endpunkte der Schneide ruht?\textsuperscript{*})\begingroup\renewcommand{\thefootnote}{}\footnotetext{\textsuperscript{*}) Denn eine Azimuthaländerung der Schwingung kann nur entweder auf diese Weise oder durch eine mit Reibung verbundene Schiebung der Schneide möglich werden. Es ist daher auch die Frage zu beantworten, oder vielmehr es ist die Hauptfrage, bei welcher Schneidenlänge die zur Schiebung nöthige Ueberwindung der Reibung möglich wird.}\endgroup\ Mir scheinen
die damit zusammenhängenden Probleme nicht ganz leicht, wenigstens ihre
Erledigung eines Zeitaufwandes bedürftig, welchen ich ihnen jetzt nicht widmen
kann.
\par\noindent .\ .\ .\ .\ .\ .\ .\ .\ .\ .\ .\ .\ .\ .\ .\ .\ .\ .\ .\ .\ .
\HistoricalFoot{}{5\textsuperscript{*}}
\clearpage\endgroup
